\documentclass[14pt,a4paper]{extarticle}
% =============================================================
% preamble.tex — оформление отчётов по СТП БГУИР 01-2024
% =============================================================

% Шрифты: под pdflatex — Times через tempora + T2A,
% под xelatex/tectonic — системный Times New Roman (запасной вариант STIX Two Text)
\usepackage{iftex}
\ifPDFTeX
  \usepackage[utf8]{inputenc}
  \usepackage[T2A]{fontenc}
  \usepackage{tempora}
\else
  \usepackage{fontspec}
  \IfFontExistsTF{Times New Roman}
    {\setmainfont{Times New Roman}}
    {\setmainfont{STIX Two Text}}
  \IfFontExistsTF{DejaVu Sans Mono}
    {\setmonofont{DejaVu Sans Mono}[Scale=0.82]}
    {\setmonofont{DejaVuSansMono}[Extension=.ttf, BoldFont=*-Bold,
                                  ItalicFont=*-Oblique, Scale=0.82]}
\fi
\usepackage[russian]{babel}

% Поля: слева 30 мм, справа 15 мм, сверху и снизу 20 мм
\usepackage{geometry}
\geometry{left=30mm, right=15mm, top=20mm, bottom=20mm}

\usepackage{setspace}
\singlespacing

\pretolerance=1000
\tolerance=2000
\emergencystretch=3em

\usepackage{indentfirst}
\setlength{\parindent}{1.25cm}

% Номер страницы — снизу справа
\usepackage{fancyhdr}
\pagestyle{fancy}
\fancyhf{}
\fancyfoot[R]{\thepage}
\renewcommand{\headrulewidth}{0pt}
\fancypagestyle{plain}{%
  \fancyhf{}%
  \fancyfoot[R]{\thepage}%
  \renewcommand{\headrulewidth}{0pt}%
}

\usepackage{amsmath,amssymb}
\usepackage{graphicx}
\usepackage{float}
\usepackage{booktabs}
\usepackage{tabularx}
\usepackage{multirow}
\usepackage{longtable}
\usepackage{array}

\usepackage{caption}
\captionsetup[figure]{labelsep=endash, justification=centering,
                      font={normalsize}, name=Рисунок}
\captionsetup[table]{labelsep=endash, justification=raggedright,
                     singlelinecheck=false, font={normalsize}, name=Таблица}

\usepackage{titlesec}
\titleformat{\section}
  {\newpage\normalfont\fontsize{14}{18}\selectfont\bfseries}
  {\thesection}{0.5em}{\MakeUppercase}
\titlespacing*{\section}{1.25cm}{0pt}{14pt}
\titleformat{\subsection}
  {\normalfont\fontsize{14}{18}\selectfont\bfseries}
  {\thesubsection}{0.5em}{}
\titlespacing*{\subsection}{1.25cm}{14pt}{14pt}
\titleformat{\subsubsection}
  {\normalfont\fontsize{14}{18}\selectfont\bfseries}
  {\thesubsubsection}{0.5em}{}
\titlespacing*{\subsubsection}{1.25cm}{14pt}{14pt}

\newcommand{\sectioncenter}[1]{%
  \newpage
  \phantomsection
  \addcontentsline{toc}{section}{#1}
  \begin{center}\textbf{\MakeUppercase{#1}}\end{center}
  \vspace{14pt}
}

\usepackage{tocloft}
\renewcommand{\cftsecfont}{\normalfont}
\renewcommand{\cftsecpagefont}{\normalfont}
\renewcommand{\cftsubsecfont}{\normalfont}
\renewcommand{\cftsubsecpagefont}{\normalfont}
\renewcommand{\cftsecleader}{\cftdotfill{\cftdotsep}}
\renewcommand{\cftdotsep}{1}
% babel переопределяет заголовки после преамбулы, поэтому правим через \addto
\addto\captionsrussian{\renewcommand{\contentsname}{\hfill\textbf{СОДЕРЖАНИЕ}\hfill}}
\renewcommand{\cftaftertoctitle}{\hfill}

\usepackage{enumitem}
\setlist[itemize]{noitemsep, topsep=0pt, leftmargin=1.25cm, label=--}
\setlist[enumerate]{noitemsep, topsep=0pt, leftmargin=1.25cm}

% Листинги кода и вывода программы
\usepackage{listings}
\usepackage{fancyvrb}
\lstset{
  basicstyle=\ttfamily\footnotesize,
  frame=single,
  breaklines=true,
  showstringspaces=false,
  inputencoding=utf8,
  extendedchars=true,
  columns=fullflexible,
  keepspaces=true,
  captionpos=t
}
\ifPDFTeX
\lstset{
  literate={а}{{\cyra}}1 {б}{{\cyrb}}1 {в}{{\cyrv}}1 {г}{{\cyrg}}1
           {д}{{\cyrd}}1 {е}{{\cyre}}1 {ж}{{\cyrzh}}1 {з}{{\cyrz}}1
           {и}{{\cyri}}1 {й}{{\cyrishrt}}1 {к}{{\cyrk}}1 {л}{{\cyrl}}1
           {м}{{\cyrm}}1 {н}{{\cyrn}}1 {о}{{\cyro}}1 {п}{{\cyrp}}1
           {р}{{\cyrr}}1 {с}{{\cyrs}}1 {т}{{\cyrt}}1 {у}{{\cyru}}1
           {ф}{{\cyrf}}1 {х}{{\cyrh}}1 {ц}{{\cyrc}}1 {ч}{{\cyrch}}1
           {ш}{{\cyrsh}}1 {щ}{{\cyrshch}}1 {ъ}{{\cyrhrdsn}}1
           {ы}{{\cyrery}}1 {ь}{{\cyrsftsn}}1 {э}{{\cyrerev}}1
           {ю}{{\cyryu}}1 {я}{{\cyrya}}1 {ё}{{\cyryo}}1
           {А}{{\CYRA}}1 {Б}{{\CYRB}}1 {В}{{\CYRV}}1 {Г}{{\CYRG}}1
           {Д}{{\CYRD}}1 {Е}{{\CYRE}}1 {Ж}{{\CYRZH}}1 {З}{{\CYRZ}}1
           {И}{{\CYRI}}1 {Й}{{\CYRISHRT}}1 {К}{{\CYRK}}1 {Л}{{\CYRL}}1
           {М}{{\CYRM}}1 {Н}{{\CYRN}}1 {О}{{\CYRO}}1 {П}{{\CYRP}}1
           {Р}{{\CYRR}}1 {С}{{\CYRS}}1 {Т}{{\CYRT}}1 {У}{{\CYRU}}1
           {Ф}{{\CYRF}}1 {Х}{{\CYRH}}1 {Ц}{{\CYRC}}1 {Ч}{{\CYRCH}}1
           {Ш}{{\CYRSH}}1 {Щ}{{\CYRSHCH}}1 {Ъ}{{\CYRHRDSN}}1
           {Ы}{{\CYRERY}}1 {Ь}{{\CYRSFTSN}}1 {Э}{{\CYREREV}}1
           {Ю}{{\CYRYU}}1 {Я}{{\CYRYA}}1 {Ё}{{\CYRYO}}1
           {—}{{---}}1 {–}{{--}}1 {…}{{\ldots}}1
}
\fi

% Подписи листингов — в том же стиле, что таблицы: «Листинг X.Y — Название»
\captionsetup[lstlisting]{labelsep=endash, justification=raggedright,
                          singlelinecheck=false, font={normalsize}}
\renewcommand{\lstlistingname}{Листинг}
\AtBeginDocument{\numberwithin{lstlisting}{section}}

\lstdefinestyle{python}{language=Python, keywordstyle=\bfseries}
\lstdefinestyle{shell}{language=bash, keywordstyle=\mdseries}
\lstdefinestyle{plain}{language={}, basicstyle=\ttfamily\scriptsize}

% Абзац после листинга должен сохранять отступ 1,25 см (СТП 2.1.1)
\makeatletter
\AtBeginDocument{\let\@endpetrue\@endpefalse}
\makeatother

% Подпись к блоку verbatim: \captionof вне плавающего окружения
% глобально обнуляет \parindent, поэтому отступ восстанавливается вручную
\newcommand{\listingcaption}[2]{%
  \captionof{lstlisting}{#1}\label{#2}%
  \setlength{\parindent}{1.25cm}%
}

\usepackage{hyperref}
\hypersetup{colorlinks=true, linkcolor=black, urlcolor=blue, citecolor=black}

\numberwithin{figure}{section}
\numberwithin{table}{section}
\numberwithin{equation}{section}

% Титульный лист отчёта по лабораторной работе
\newcommand{\makelabtitle}[2]{%
  \begin{titlepage}
  \thispagestyle{empty}
  \begin{center}

  МИНИСТЕРСТВО ОБРАЗОВАНИЯ РЕСПУБЛИКИ БЕЛАРУСЬ

  \vspace{0.5em}
  Учреждение образования

  \vspace{0.5em}
  \textbf{<<Белорусский государственный университет информатики и радиоэлектроники>>}

  \vspace{2em}
  Факультет компьютерных систем и сетей

  \vspace{0.5em}
  Кафедра программного обеспечения информационных технологий
  \end{center}

  \vspace{3em}
  \begin{center}
  \textbf{ОТЧЕТ}

  \vspace{0.5em}
  по лабораторной работе №#1

  \vspace{0.5em}
  по дисциплине

  \vspace{0.5em}
  <<Информационные и технические средства защиты авторских прав>>

  \vspace{0.5em}
  <<#2>>
  \end{center}

  \vspace{5em}
  \begin{flushright}
  \begin{minipage}{0.5\textwidth}
  \textbf{Выполнил:}\\
  Лебедевич Артем Владимирович\\
  магистрант группы 556301

  \vspace{0.6em}
  \textbf{Преподаватель:}\\
  Ярмолик Вячеслав Николаевич\\
  доктор технических наук\\
  профессор кафедры ПОИТ
  \end{minipage}
  \end{flushright}

  \vfill
  \begin{center}Минск, 2026~г.\end{center}
  \end{titlepage}
}

% Титульный лист реферата
\newcommand{\makerefattitle}[1]{%
  \begin{titlepage}
  \thispagestyle{empty}
  \begin{center}

  МИНИСТЕРСТВО ОБРАЗОВАНИЯ РЕСПУБЛИКИ БЕЛАРУСЬ

  \vspace{0.5em}
  Учреждение образования

  \vspace{0.5em}
  \textbf{<<Белорусский государственный университет информатики и радиоэлектроники>>}

  \vspace{2em}
  Факультет компьютерных систем и сетей

  \vspace{0.5em}
  Кафедра программного обеспечения информационных технологий
  \end{center}

  \vspace{3em}
  \begin{center}
  \textbf{РЕФЕРАТ}

  \vspace{0.5em}
  по дисциплине

  \vspace{0.5em}
  <<Информационные и технические средства защиты авторских прав>>

  \vspace{1.5em}
  на тему

  \vspace{0.5em}
  \textbf{<<#1>>}
  \end{center}

  \vspace{5em}
  \begin{flushright}
  \begin{minipage}{0.5\textwidth}
  \textbf{Выполнил:}\\
  Лебедевич Артем Владимирович\\
  магистрант группы 556301

  \vspace{0.6em}
  \textbf{Преподаватель:}\\
  Ярмолик Вячеслав Николаевич\\
  доктор технических наук\\
  профессор кафедры ПОИТ
  \end{minipage}
  \end{flushright}

  \vfill
  \begin{center}Минск, 2026~г.\end{center}
  \end{titlepage}
}


\begin{document}

\makelabtitle{2}{Депозитные операции с физическими лицами}

\tableofcontents

\section{Цель работы}

Реализовать микросервис депозитных операций с физическими лицами: справочник <<План
счетов>>, форму заключения депозитного договора с контролем вводимых данных, открытие
счетов клиента по правилам банковского учёта, проводки по схеме <<Депозитная программа>>,
счёт фонда развития банка и процедуру <<Закрытие банковского дня>> с отчётом о состоянии
счетов. Данные о клиенте берутся из модуля <<Клиенты>> первой лабораторной работы.

\section{Задание}

Вариант №\,4 соответствует банку <<Решение>> в списке банков Республики Беларусь
(\url{https://select.by/banki}). Для него выбраны две депозитные программы для физических
лиц (таблица~\ref{tab:products}); условия приближены к действующим предложениям банка.

\begin{table}[H]
\caption{Депозитные программы банка <<Решение>>}
\label{tab:products}
\centering
\small
\begin{tabularx}{\textwidth}{lXX}
\toprule
 & <<Свободное накопление>> & <<Online-решение>> \\
\midrule
Вид вклада & отзывный, до востребования & срочный безотзывный \\
Выплата процентов & ежемесячно & в конце срока \\
Досрочный отзыв & возможен в любой день & невозможен \\
Ставка, BYN & 7,00\,\% годовых & 12,90\,\% годовых \\
Ставка, USD & 1,50\,\% годовых & 3,00\,\% годовых \\
Минимальная сумма & 100 & 500 \\
Срок договора & 1--12 месяцев & 3--13 месяцев \\
Счёт основной суммы & 3404 & 3414 \\
Счёт процентов & 3470 & 3471 \\
\bottomrule
\end{tabularx}
\end{table}

Требования задания:
\begin{enumerate}
\item разработать справочник <<План счетов>> и внести в него счета, необходимые для
обслуживания выбранных депозитных линий;
\item разработать форму заключения депозитного договора (вид депозита и валюта из
справочников, номер договора, даты начала и окончания, срок, сумма, процент) с контролем
корректности данных;
\item при заключении договора создавать не менее двух счетов --- для основной суммы и
для процентов --- с 13-значным номером, кодом и активностью из плана счетов, дебетом,
кредитом, сальдо и названием по ФИО клиента;
\item предусмотреть счёт фонда развития банка (СФРБ) со стартовым капиталом 1~млн
белорусских рублей;
\item разработать процедуру <<Закрытие банковского дня>> и отчёт о состоянии счетов; при
активных депозитных программах в нём должно быть не менее шести счетов.
\end{enumerate}

\section{Проектирование}

\subsection{Декомпозиция на микросервисы}

Депозитные операции затрагивают три разных области ответственности, и каждой
соответствует свой микросервис со своей базой данных (рисунок~\ref{fig:services}):
\begin{itemize}
\item \texttt{client-service} --- клиенты банка (лабораторная работа №\,1);
\item \texttt{account-service} --- главная книга: план счетов, лицевые счета, проводки,
банковский день;
\item \texttt{deposit-service} --- депозитные программы и договоры, расчёт процентов.
\end{itemize}

\begin{figure}[H]
\centering
\begin{tikzpicture}[node distance=10mm and 16mm]
\node[svc] (dep) at (0, 0) {deposit-service\\[-1pt]{\footnotesize договоры, проценты}\\[-3pt]{\footnotesize порт 8083}};
\node[svc] (cli) at (-5.9, 0) {client-service\\[-1pt]{\footnotesize клиенты (ЛР1)}\\[-3pt]{\footnotesize порт 8081}};
\node[svc] (acc) at (5.9, 0) {account-service\\[-1pt]{\footnotesize счета, проводки, день}\\[-3pt]{\footnotesize порт 8082}};
\node[db, below=7mm of cli] (dbc) {H2 clients};
\node[db, below=7mm of dep] (dbd) {H2 deposits};
\node[db, below=7mm of acc] (dba) {H2 accounts};
\node[box, above=11mm of dep, fill=codeback, minimum width=3.3cm] (ui) {Web-клиент депозитов\\[-1pt]{\footnotesize Vue 3}};
\draw[flow] (cli) -- (dbc);
\draw[flow] (dep) -- (dbd);
\draw[flow] (acc) -- (dba);
\draw[flow] (ui) -- node[right, note] {REST} (dep);
\draw[flow] (dep) -- node[above, note] {клиент} (cli);
\draw[flow] ([yshift=2.5mm]dep.east) -- node[above, note] {счета, проводки} ([yshift=2.5mm]acc.west);
\draw[flow, dashed] ([yshift=-2.5mm]acc.west) -- node[below, note] {новый день} ([yshift=-2.5mm]dep.east);
\end{tikzpicture}
\caption{Микросервисы депозитного контура и REST-вызовы между ними}
\label{fig:services}
\end{figure}

Главная книга выделена в отдельный сервис намеренно. Со счетами работают и депозиты, и
кредиты следующей лабораторной работы, а правила учёта --- расчёт сальдо по активности
счёта, запрет отрицательного остатка, неделимость проводки --- должны соблюдаться в одном
месте независимо от того, какой сервис инициировал операцию. Сервис депозитов не хранит ни
оборотов, ни остатков: он знает только номера счетов своего договора и просит главную книгу
провести операцию.

Сервисы общаются синхронно по REST. Клиенты смежных сервисов описаны декларативными
интерфейсами Spring (\texttt{@HttpExchange}) в общей библиотеке \texttt{bank-common}; там же
находится единый формат ошибки, который позволяет передать сообщение главной книги через
сервис депозитов до web-клиента без искажений.

\subsection{План счетов и номер счёта}

Справочник <<План счетов>> (таблица~\ref{tab:chart}) хранится в базе сервиса счетов и
заполняется при старте. Активность балансового счёта определяет правило расчёта сальдо:
у активного счёта остаток равен дебету минус кредит, у пассивного --- кредиту минус дебет.

\begin{table}[H]
\caption{Счета плана счетов, используемые депозитными программами}
\label{tab:chart}
\centering
\small
\begin{tabularx}{\textwidth}{clXl}
\toprule
Код & Активность & Наименование & Назначение в работе \\
\midrule
1010 & активный & Денежные средства в кассе & касса банка \\
3404 & пассивный & Вклады (депозиты) до востребования физических лиц & сумма отзывного вклада \\
3414 & пассивный & Срочные вклады (депозиты) физических лиц & сумма срочного вклада \\
3470 & пассивный & Начисленные процентные расходы по вкладам до востребования & проценты отзывного вклада \\
3471 & пассивный & Начисленные процентные расходы по срочным вкладам & проценты срочного вклада \\
7327 & пассивный & Фонд развития банка & СФРБ \\
\bottomrule
\end{tabularx}
\end{table}

Номер лицевого счёта состоит из 13 цифр:
\begin{equation}
\underbrace{b_1 b_2 b_3 b_4}_{\text{балансовый счёт}}\;
\underbrace{c_1 c_2 c_3 c_4 c_5}_{\text{код клиента}}\;
\underbrace{n_1 n_2 n_3}_{\text{номер счёта клиента}}\;
\underbrace{k}_{\text{ключ}}, \qquad
k = \Bigl(3 \sum_{i=1}^{12} d_i w_i\Bigr) \bmod 10,
\end{equation}
где $d_i$~--- первые 12 цифр номера, $w_i$~--- веса, циклически повторяющие
последовательность $7, 1, 3$. Код клиента равен его идентификатору в модуле <<Клиенты>>, у
собственных счетов банка он нулевой. Веса 7, 1, 3 и множитель 3 взаимно просты с 10,
поэтому ошибка в любой одной цифре номера меняет ключ и обнаруживается при проверке.

\subsection{Схема данных}

Базы сервисов независимы (рисунок~\ref{fig:er}). Договор ссылается на клиента и на счета
не внешними ключами, а идентификатором клиента и номерами счетов: целостность таких ссылок
обеспечивается проверкой через REST в момент заключения договора.

\begin{figure}[H]
\centering
\begin{tikzpicture}
\node[note, anchor=west] at (0, 0.45) {\textit{account-service}};
\node[entity, text width=3.6cm, anchor=north west] (chart) at (0, 0) {\textbf{chart\_of\_accounts}\nodepart{two} code \hfill PK\\ name\\ activity};
\node[entity, text width=4.3cm, anchor=north west] (account) at (5.0, 0) {\textbf{account}\nodepart{two} id \hfill PK\\ number \hfill UQ\\ chart\_code \hfill FK\\ currency\_code \hfill FK\\ name, contract\_ref\\ owner\_code, owner\_seq \hfill UQ\\ debit, credit, opened\_on};
\node[entity, text width=3.0cm, anchor=north west] (currency) at (10.7, 0) {\textbf{currency}\nodepart{two} code \hfill PK\\ name};
\node[entity, text width=3.6cm, anchor=north west] (entry) at (0, -4.3) {\textbf{operation\_entry}\nodepart{two} id \hfill PK\\ operation\_id \hfill FK\\ account\_id \hfill FK\\ side, amount};
\node[entity, text width=4.3cm, anchor=north west] (operation) at (5.0, -4.3) {\textbf{operation}\nodepart{two} id \hfill PK\\ batch\_key \hfill FK\\ bank\_date\\ description, contract\_ref};
\node[entity, text width=3.0cm, anchor=north west] (batch) at (10.7, -4.3) {\textbf{posting\_batch}\nodepart{two} batch\_key \hfill PK\\ created\_at};
\node[entity, text width=3.0cm, anchor=north west] (day) at (10.7, -1.9) {\textbf{bank\_day}\nodepart{two} id \hfill PK\\ bank\_date};
\draw[flow] (account.west |- chart.east) -- (chart.east);
\draw[flow] (account.east |- currency.west) -- (currency.west);
\draw[flow] (entry.east |- operation.west) -- (operation.west);
\draw[flow] (entry.north) -- ++(0, 0.55) -| ([xshift=-12mm]account.south);
\draw[flow] (operation.east |- batch.west) -- (batch.west);

\node[note, anchor=west] at (0, -7.35) {\textit{deposit-service}};
\node[entity, text width=4.4cm, anchor=north west] (product) at (0, -7.8) {\textbf{deposit\_product}\nodepart{two} id \hfill PK\\ name, kind, currency\\ rate, min\_amount\\ min\_term\_months\\ max\_term\_months};
\node[entity, text width=7.2cm, anchor=north west] (contract) at (6.5, -7.8) {\textbf{deposit\_contract}\nodepart{two} id \hfill PK\\ number \hfill UQ\\ product\_id \hfill FK\\ client\_id, client\_name \hfill $\to$ client-service\\ currency, amount, rate, term\_months\\ start\_date, end\_date, status\\ main\_account, interest\_account \hfill $\to$ account\\ accrued, paid, paid\_months};
\draw[flow] (contract.west |- product.east) -- (product.east);
\end{tikzpicture}
\caption{Схемы данных сервисов счетов и депозитов}
\label{fig:er}
\end{figure}

\subsection{Схема проводок}

Проводки депозитной программы выполняются в точности по упрощённой схеме из задания
(таблица~\ref{tab:postings}). Знак <<$+$>> означает увеличение остатка счёта, <<$-$>> ---
уменьшение: для активной кассы остаток растёт по дебету, для трёх пассивных счетов --- по
кредиту. Касса в схеме работает как транзитный счёт: всё, что внесено, в тот же момент
зачисляется на счёт клиента или выдаётся, поэтому её остаток возвращается к нулю.

\begin{table}[H]
\caption{Проводки процесса <<Депозитная программа>>}
\label{tab:postings}
\centering
\footnotesize
\setlength{\tabcolsep}{4.5pt}
\begin{tabular}{clcccccccc}
\toprule
 & & \multicolumn{2}{c}{Касса (А)} & \multicolumn{2}{c}{СФРБ (П)} & \multicolumn{2}{c}{Счёт вклада (П)} & \multicolumn{2}{c}{Проценты (П)} \\
\cmidrule(lr){3-4}\cmidrule(lr){5-6}\cmidrule(lr){7-8}\cmidrule(lr){9-10}
№ & Операция & Дт & Кт & Дт & Кт & Дт & Кт & Дт & Кт \\
\midrule
1 & Внесение денег в кассу & $+$ & & & & & & & \\
2 & Перевод денег с кассы на текущий счёт & & $-$ & & & & $+$ & & \\
3 & Использование денег банком & & & & $+$ & $-$ & & & \\
\midrule
4 & Начисление процентов по депозиту & & & $-$ & & & & & $+$ \\
5 & Перевод процентов в кассу & $+$ & & & & & & $-$ & \\
6 & Вывод процентов из кассы & & $-$ & & & & & & \\
\midrule
7 & Окончание депозита & & & $-$ & & & $+$ & & \\
8 & Перевод депозита в кассу & $+$ & & & & $-$ & & & \\
9 & Вывод денег из кассы & & $-$ & & & & & & \\
\bottomrule
\end{tabular}
\end{table}

Операции 1--3 выполняются при заключении договора, 4~--- при каждом закрытии дня,
5--6~--- в день выплаты процентов, 7--9~--- по окончании срока или при досрочном отзыве.

\subsection{Расчёт процентов}

Проценты простые, начисляются за фактическое число дней при 365 днях в году. К дате
$t$ по договору на сумму $S$ со ставкой $r$ процентов годовых, заключённому в день $t_0$,
накоплено
\begin{equation}
I(t) = \operatorname{round}\Bigl(\frac{S \cdot r \cdot (t - t_0)}{100 \cdot 365},\ 2\Bigr),
\label{eq:interest}
\end{equation}
а проводка очередного дня делается на разность $I(t)$ и уже начисленной суммы. Такой
накопительный расчёт устойчивее суммирования дневных процентов: копейки округления не
теряются, а пропущенный или повторно обработанный день не искажает итог. Например, вклад
10\,000 BYN под 7\,\% за 31 день приносит $10000 \cdot 7 \cdot 31 / 36500 = 59{,}45$ BYN.

По отзывному вкладу накопленные проценты выплачиваются, когда истекает очередной полный
месяц договора; по срочному --- только в день окончания, вместе с возвратом суммы вклада.

\subsection{Закрытие банковского дня}

Банковский день хранится в сервисе счетов, но сам этот сервис не знает правил
депозитов. Поэтому закрытие дня устроено как оповещение
(рисунок~\ref{fig:dayclose}): сервис счетов объявляет участникам дату нового дня, каждый
участник начисляет проценты и проводит выплаты через обычный REST главной книги и
возвращает протокол сделанного, после чего дата сдвигается.

\begin{figure}[H]
\centering
\begin{tikzpicture}[font=\footnotesize]
\foreach \x/\name in {0/Web-клиент, 5.6/account-service, 11.2/deposit-service} {
  \node[box, minimum width=2.9cm] at (\x, 0) {\name};
  \draw[black!50, dashed] (\x, -0.4) -- (\x, -6.6);
}
\draw[flow] (0, -1.1) -- node[above] {POST /api/bank-day/close} (5.6, -1.1);
\draw[flow] (5.6, -2.0) -- node[above] {POST /api/internal/day-close} (11.2, -2.0);
\draw[flow] (11.2, -3.0) -- node[above] {POST /api/operations} (5.6, -3.0);
\draw[flow, dashed] (5.6, -3.7) -- node[above] {проводки проведены} (11.2, -3.7);
\draw[flow, dashed] (11.2, -4.6) -- node[above] {протокол операций дня} (5.6, -4.6);
\node[note, anchor=west] at (5.75, -5.3) {bank\_date := date};
\draw[flow, dashed] (5.6, -6.1) -- node[above] {новая дата и протокол} (0, -6.1);
\end{tikzpicture}
\caption{Взаимодействие сервисов при закрытии банковского дня}
\label{fig:dayclose}
\end{figure}

Распределённая процедура должна переживать сбои, поэтому в неё заложены три свойства.
\begin{enumerate}
\item \textbf{Идемпотентность.} Каждый пакет проводок несёт ключ вида
\texttt{DEP-Д-000001-ACCRUE-2026-10-02}. Главная книга запоминает ключи и на повторный запрос
возвращает уже проведённые операции, не проводя их второй раз.
\item \textbf{Дата сдвигается последней.} Если участник ответил ошибкой, день остаётся
открытым и закрытие можно повторить: уже обработанные договоры повторно не начисляются.
\item \textbf{Навёрстывание.} Недоступный участник не блокирует закрытие дня. Его расчёт
накопительный (формула~\ref{eq:interest}), поэтому при следующем закрытии он начислит
проценты сразу за все пропущенные дни.
\end{enumerate}

\subsection{REST-интерфейсы}

Операции двух сервисов приведены в таблице~\ref{tab:rest}.

\begin{table}[H]
\caption{REST-интерфейсы сервисов счетов и депозитов}
\label{tab:rest}
\centering
\footnotesize
\begin{tabularx}{\textwidth}{llX}
\toprule
Метод & Адрес & Назначение \\
\midrule
\multicolumn{3}{l}{\textit{account-service}} \\
GET & \texttt{/api/chart-of-accounts}, \texttt{/api/currencies} & справочники <<План счетов>> и валют \\
GET & \texttt{/api/accounts?contractRef=} & счета договора или клиента \\
GET & \texttt{/api/accounts/system} & счёт кассы или фонда развития в заданной валюте \\
POST & \texttt{/api/accounts} & открытие лицевого счёта \\
POST & \texttt{/api/operations} & атомарный пакет проводок с ключом идемпотентности \\
GET & \texttt{/api/operations} & журнал проводок \\
GET & \texttt{/api/report} & отчёт о состоянии счетов с итогами оборотов \\
GET & \texttt{/api/bank-day} & текущий банковский день \\
POST & \texttt{/api/bank-day/close?days=} & закрытие одного или нескольких дней \\
\midrule
\multicolumn{3}{l}{\textit{deposit-service}} \\
GET & \texttt{/api/products}, \texttt{/api/meta} & виды депозитов, валюты, банковский день, номер договора \\
GET & \texttt{/api/clients} & список клиентов из модуля <<Клиенты>> \\
GET & \texttt{/api/deposits}, \texttt{/api/deposits/\{id\}} & договоры; договор со счетами и проводками \\
POST & \texttt{/api/deposits} & заключение договора \\
POST & \texttt{/api/deposits/\{id\}/withdraw} & досрочный отзыв вклада \\
POST & \texttt{/api/internal/day-close} & обработка нового банковского дня \\
\bottomrule
\end{tabularx}
\end{table}

\section{Реализация}

\subsection{Структура проекта}

К проекту первой работы добавлены три модуля Maven (листинг~\ref{lst:tree}).

\listingcaption{Модули лабораторной работы №\,2}{lst:tree}
\begin{Verbatim}[fontsize=\footnotesize, frame=single, rulecolor=\color{codeframe}]
bank-common/                общая библиотека сервисов
  api/                      ApiError, BankException, обработчик ошибок REST
  client/, ledger/          REST-клиенты ClientApi и LedgerApi, записи запросов и ответов
account-service/            главная книга
  domain/                   Account, ChartAccount, Operation, AccountNumbers, Activity
  service/                  LedgerService, BankDayService, DayCloseNotifier, LedgerInitializer
  web/                      LedgerController, BankDayController
  resources/                schema.sql, data.sql (план счетов), static/ (отчёт по счетам)
deposit-service/            депозиты
  domain/                   DepositProduct, DepositContract, DepositKind
  service/                  DepositService, DepositPostings, InterestCalculator, DayCloseProcessor
  web/                      DepositController, ReferenceController, DayCloseController
  resources/                schema.sql, data.sql (программы банка), static/ (форма договора)
\end{Verbatim}

\subsection{Главная книга}

Операция в главной книге --- это набор строк <<счёт, сторона, сумма>>. Сервис принимает не
одиночную операцию, а пакет (листинг~\ref{lst:post}): заключение договора состоит из трёх
операций схемы, и они должны быть проведены все вместе либо не проведены вовсе. Пакет
выполняется в одной транзакции базы данных. Счета выбираются запросом
\texttt{select ... for update}, чтобы два параллельных пакета не потеряли обновление
оборотов. После каждой операции проверяется остаток затронутых счетов: отрицательное сальдо
означает нехватку средств, и тогда исключение откатывает весь пакет вместе с записью о ключе.

\begin{lstlisting}[style=java, caption={Проведение пакета операций (\texttt{LedgerService.post})}, label={lst:post}]
@Transactional
public List<OperationView> post(BatchRequest request) {
    if (batches.existsById(request.batchKey())) {          // повтор запроса: операции уже проведены
        return operations.findByBatchKeyOrderById(request.batchKey()).stream().map(OperationView::from).toList();
    }
    batches.saveAndFlush(new PostingBatch(request.batchKey()));

    LocalDate date = request.bankDate() != null ? request.bankDate() : bankDay.current();
    Map<String, Account> locked = new LinkedHashMap<>();
    List<Operation> posted = new ArrayList<>();
    for (OperationRequest item : request.operations()) {
        Operation operation = new Operation();
        ...
        Set<Account> touched = new LinkedHashSet<>();
        for (EntryRequest entry : item.entries()) {
            Account account = locked.computeIfAbsent(entry.account(), number -> accounts.findByNumberForUpdate(number)
                    .orElseThrow(() -> BankException.notFound("Счёт " + number + " не найден")));
            ...
            account.apply(entry.side(), entry.amount());
            operation.addEntry(account, entry.side(), entry.amount());
            touched.add(account);
        }
        touched.forEach(account -> checkNotOverdrawn(account, item.description()));
        posted.add(operations.save(operation));
    }
    operations.flush();
    return posted.stream().map(OperationView::from).toList();
}
\end{lstlisting}

При старте сервис открывает в каждой валюте кассу и фонд развития банка и зачисляет в
фонд стартовый капитал. Размер капитала задаётся параметром \texttt{bank.fund.capital}; для
этой работы сервис запускается со значением 1\,000\,000 BYN.

\subsection{Заключение депозитного договора}

Сервис депозитов проверяет условия договора, запрашивает клиента у модуля <<Клиенты>>,
открывает в главной книге два счёта и проводит операции 1--3 схемы
(листинг~\ref{lst:open}). Балансовые счета берутся из вида вклада: 3404 и 3470 для вклада
до востребования, 3414 и 3471 для срочного. Названием счёта служит ФИО клиента.

\begin{lstlisting}[style=java, caption={Открытие счетов и проводки приёма вклада (\texttt{DepositService.open})}, label={lst:open}]
ClientInfo client = findClient(request.clientId());

DepositKind kind = product.getKind();
String number = request.number();
String currency = product.getCurrency();
String main = ledger.open(new OpenAccount(kind.getMainChartCode(), currency,
        client.id().intValue(), client.fullName(), number)).number();
String interest = ledger.open(new OpenAccount(kind.getInterestChartCode(), currency,
        client.id().intValue(), client.fullName(), number)).number();
post(key(number, "OPEN"), number, today,
        DepositPostings.open(cash(currency), fund(currency), main, request.amount()));
\end{lstlisting}

Сами проводки собраны в классе \texttt{DepositPostings}, который читается как таблица
схемы (листинг~\ref{lst:postings}).

\begin{lstlisting}[style=java, caption={Проводки приёма вклада и выплаты процентов (\texttt{DepositPostings})}, label={lst:postings}]
public static List<Posting> open(String cash, String fund, String main, BigDecimal amount) {
    return List.of(
            Posting.of("Внесение денег в кассу", Entry.debit(cash, amount)),
            Posting.of("Перевод денег с кассы на текущий счёт", Entry.credit(cash, amount), Entry.credit(main, amount)),
            Posting.of("Использование денег банком", Entry.debit(main, amount), Entry.credit(fund, amount)));
}

public static Posting accrual(String fund, String interest, BigDecimal amount) {
    return Posting.of("Начисление процентов по депозиту", Entry.debit(fund, amount), Entry.credit(interest, amount));
}

public static List<Posting> payout(String cash, String interest, BigDecimal amount) {
    return List.of(
            Posting.of("Перевод процентов в кассу", Entry.debit(interest, amount), Entry.debit(cash, amount)),
            Posting.of("Вывод процентов из кассы", Entry.credit(cash, amount)));
}
\end{lstlisting}

Открытие счёта в главной книге идемпотентно по паре <<договор, балансовый счёт>>, а пакет
проводок --- по ключу \texttt{DEP-<номер>-OPEN}. Если сбой произойдёт после обращения к
главной книге, но до сохранения договора, повтор запроса с тем же номером договора не
создаст ни вторых счетов, ни вторых проводок.

\subsection{Контроль корректности данных договора}

Как и в первой работе, проверка двойная. Web-клиент проверяет форму до отправки, сервис
повторяет проверки и добавляет те, что возможны только на сервере
(таблица~\ref{tab:checks}). Ставка и дата начала в форме недоступны для редактирования,
но сервис всё равно сверяет их: запрос может прийти мимо формы или из устаревшей страницы.

\begin{table}[H]
\caption{Проверки формы депозитного договора}
\label{tab:checks}
\centering
\small
\begin{tabularx}{\textwidth}{lXc}
\toprule
Поле & Правило & Где проверяется \\
\midrule
Клиент & выбран; существует в модуле <<Клиенты>> & клиент, сервис \\
Номер договора & формат \texttt{Д-000001}; не занят & клиент, сервис \\
Вид депозита & запись справочника; доступен в выбранной валюте & клиент, сервис \\
Сумма & денежная; не меньше минимума программы & клиент, сервис \\
Срок & целое число месяцев в пределах программы & клиент, сервис \\
Процент & равен ставке программы & сервис \\
Дата начала & равна текущему банковскому дню & сервис \\
Дата окончания & равна дате начала плюс срок & сервис \\
\bottomrule
\end{tabularx}
\end{table}

\subsection{Обработка банковского дня}

Обработка договора на дату нового дня показана в листинге~\ref{lst:day}. Каждый договор
обрабатывается в собственной транзакции, а ключи пакетов включают дату, поэтому повторный
вызов для того же дня ничего не меняет.

\begin{lstlisting}[style=java, caption={Начисление, выплата и закрытие договора (\texttt{DepositService.processDay})}, label={lst:day}]
boolean matured = !date.isBefore(contract.getEndDate());
LocalDate upTo = matured ? contract.getEndDate() : date;

BigDecimal target = InterestCalculator.accrued(contract.getAmount(), contract.getRate(), contract.getStartDate(), upTo);
BigDecimal delta = target.subtract(contract.getAccrued());
if (delta.signum() > 0) {
    post(key(number, "ACCRUE-" + date), number, date,
            List.of(DepositPostings.accrual(fund(currency), contract.getInterestAccount(), delta)));
    contract.setAccrued(target);
}

int months = InterestCalculator.fullMonths(contract.getStartDate(), upTo);
boolean monthlyPayout = contract.getProduct().getKind() == DepositKind.REVOCABLE && months > contract.getPaidMonths();
if (matured || monthlyPayout) {
    payInterest(contract, date);
    contract.setPaidMonths(months);
}
if (matured) {
    close(contract, date);
}
\end{lstlisting}

\section{Тестирование}

Сервисы покрыты 69 тестами (таблица~\ref{tab:tests}). Чтобы тестировать сервис депозитов
без запуска соседних микросервисов, смежные сервисы заменяются двойниками: клиент модуля
<<Клиенты>> --- заглушкой Mockito, а главная книга --- классом \texttt{InMemoryLedger},
который реализует тот же интерфейс \texttt{LedgerApi} в памяти и соблюдает те же правила
(сальдо по активности, запрет отрицательного остатка, идемпотентность). Это позволяет
сверять в тестах обороты счетов с таблицей~\ref{tab:postings}.

\begin{table}[H]
\caption{Тесты сервисов счетов и депозитов}
\label{tab:tests}
\centering
\small
\begin{tabularx}{\textwidth}{lcX}
\toprule
Класс тестов & Тестов & Что проверяется \\
\midrule
\multicolumn{3}{l}{\textit{account-service}} \\
\texttt{AccountNumbersTest} & 7 & структура номера, контрольный ключ, сальдо по активности \\
\texttt{LedgerServiceIT} & 8 & счета, обороты, атомарность пакета, идемпотентность \\
\texttt{BankDayIT} & 6 & закрытие дня, недоступный и отказавший участник, REST \\
\midrule
\multicolumn{3}{l}{\textit{deposit-service}} \\
\texttt{InterestCalculatorTest} & 13 & формула процентов, полные месяцы \\
\texttt{DepositFlowIT} & 7 & проводки по схеме, выплаты, окончание и отзыв вклада \\
\texttt{DepositApiIT} & 28 & контроль данных договора, ошибки смежных сервисов \\
\midrule
Всего & 69 & \\
\bottomrule
\end{tabularx}
\end{table}

Пример теста, который проходит весь срок срочного вклада и сверяет обороты со схемой,
приведён в листинге~\ref{lst:test}; итог прогона --- в листинге~\ref{lst:run}.

\begin{lstlisting}[style=java, caption={Тест срочного вклада (\texttt{DepositFlowIT})}, label={lst:test}]
@Test
@DisplayName("Срочный вклад: проценты выплачиваются в конце срока вместе с возвратом суммы")
void irrevocableDepositPaysAtMaturity() {
    DepositView deposit = open("Д-000002", IRREVOCABLE_BYN, "5000.00", 3);

    closeDays(91);                                           // 31.12: до окончания срока один день
    assertThat(view(deposit).paid()).isEqualByComparingTo("0");
    assertThat(ledger.balance(deposit.interestAccount())).isEqualByComparingTo("160.81");

    List<String> events = closeDays(1);                      // 01.01.2027: окончание договора, 92 дня
    DepositView closed = view(deposit);
    assertThat(closed.status()).isEqualTo("CLOSED");
    assertThat(closed.paid()).isEqualByComparingTo("162.58");

    // все счета клиента обнулились, касса транзитом пропустила вклад и проценты, фонд потратил проценты
    assertTurnover(deposit.mainAccount(), "10000.00", "10000.00", "0.00");
    assertTurnover(deposit.interestAccount(), "162.58", "162.58", "0.00");
    assertTurnover(cash, "10162.58", "10162.58", "0.00");
    assertThat(ledger.balance(fund)).isEqualByComparingTo(InMemoryLedger.CAPITAL.subtract(new BigDecimal("162.58")));
}
\end{lstlisting}

\listingcaption{Результат выполнения тестов}{lst:run}
\VerbatimInput[fontsize=\fontsize{7}{8.8}\selectfont, frame=single, rulecolor=\color{codeframe}, samepage=true]{listings/tests.txt}

\section{Результаты работы}

Сценарий демонстрации: два клиента открывают вклады по разным программам, банк
работает три месяца, оба договора завершаются. Сервис счетов запущен со стартовым капиталом
фонда развития 1\,000\,000 BYN.

После старта в отчёте два собственных счёта банка в белорусских рублях --- касса и фонд
развития (рисунок~\ref{fig:start}); справочник <<План счетов>> показан на
рисунке~\ref{fig:chart}.

\screenshot{img/accounts_start.png}{Отчёт по счетам после запуска: касса и фонд развития}{fig:start}

\screenshot{img/chart.png}{Справочник <<План счетов>>}{fig:chart}

Форма заключения договора приведена на рисунке~\ref{fig:form}. Клиент выбирается из
списка, полученного от модуля <<Клиенты>>, вид депозита и валюта --- из справочников;
ставка, дата начала и дата окончания подставляются автоматически. На
рисунке~\ref{fig:errors} показан контроль данных: номер договора не соответствует формату,
сумма меньше минимальной, срок больше допустимого по программе.

\screenshot{img/form.png}{Форма заключения депозитного договора}{fig:form}

\screenshot{img/form_errors.png}{Контроль корректности данных договора}{fig:errors}

После заключения договора открывается его карточка (рисунок~\ref{fig:open}): созданы
счета 3404 и 3470 на имя клиента, проведены три операции приёма вклада. У счёта вклада
обороты по дебету и кредиту равны 10\,000, сальдо нулевое --- деньги переданы в фонд
развития, как и предписывает схема.

\screenshot{img/details_open.png}{Карточка договора: счета и проводки приёма вклада}{fig:open}

Вторым договором открыт срочный вклад <<Online-решение>> на 5\,000 BYN. В отчёте по
счетам (рисунок~\ref{fig:report-open}) теперь шесть счетов: касса, фонд развития и по два
счёта на каждый договор. Фонд развития вырос на сумму обоих вкладов до 1\,015\,000 BYN.

\screenshot{img/report_open.png}{Отчёт по счетам при двух активных депозитных программах}{fig:report-open}

Процедура <<Закрытие банковского дня>> выполнена 31 раз подряд
(рисунок~\ref{fig:report-month}). По отзывному вкладу проценты за месяц в размере 59,45 BYN
начислены и выплачены, процентный счёт обнулился. По срочному вкладу начислено
$5000 \cdot 12{,}9 \cdot 31 / 36500 = 54{,}78$ BYN, и эта сумма остаётся на счёте 3471 до
конца срока. Фонд развития уменьшился на сумму начисленных процентов.

\screenshot{img/report_month.png}{Отчёт по счетам через месяц и протокол закрытия дня}{fig:report-month}

Карточка отзывного вклада через месяц и журнал проводок главной книги показаны на
рисунках~\ref{fig:details-month} и~\ref{fig:journal}: видны ежедневные начисления, перевод
процентов в кассу и их выдача.

\screenshot[0.95\textwidth]{img/details_month.png}{Карточка отзывного вклада после выплаты процентов}{fig:details-month}

\screenshot[0.95\textwidth]{img/journal.png}{Журнал проводок главной книги}{fig:journal}

После закрытия ещё 61 дня наступает 1 января 2027 года --- дата окончания обоих
договоров. Вклады возвращены, договоры закрыты (рисунок~\ref{fig:closed}). По отзывному
вкладу за 92 дня выплачено 176,44 BYN, по срочному --- 162,58 BYN одной суммой. В итоговом
отчёте (рисунок~\ref{fig:report-closed}) сальдо всех клиентских счетов и кассы нулевое, а
фонд развития равен стартовому капиталу за вычетом выплаченных процентов:
$1\,000\,000 - 176{,}44 - 162{,}58 = 999\,660{,}98$ BYN.

\screenshot{img/list_closed.png}{Список договоров после окончания срока}{fig:closed}

\screenshot{img/report_closed.png}{Итоговый отчёт по счетам}{fig:report-closed}

\sectioncenter{Заключение}

В работе реализованы два микросервиса: главная книга банка и сервис депозитных операций.
Главная книга ведёт справочник <<План счетов>>, открывает лицевые счета с 13-значным
номером и контрольным ключом, проводит операции атомарными пакетами и формирует отчёт о
состоянии счетов. Сервис депозитов поддерживает две программы банка <<Решение>> ---
отзывный вклад с ежемесячной выплатой процентов и срочный безотзывный вклад с выплатой
в конце срока; при заключении договора он открывает два счёта клиента и выполняет проводки
строго по схеме из задания.

Процедура <<Закрытие банковского дня>> построена как взаимодействие сервисов: главная
книга объявляет новую дату, сервис депозитов начисляет и выплачивает проценты. Идемпотентные
ключи проводок, сдвиг даты после успешной обработки и накопительный расчёт процентов делают
процедуру устойчивой к повторам и к временной недоступности участника --- это те свойства,
без которых распределённая система не может надёжно работать с деньгами.

Данные о клиентах берутся из модуля первой работы по REST. Корректность проверена 69
автоматическими тестами, в том числе сверкой оборотов всех счетов со схемой проводок на
полном жизненном цикле договора.

\sectioncenter{Список использованных источников}

\begin{enumerate}
\item Об установлении Плана счетов бухгалтерского учёта в банках и небанковских
кредитно-финансовых организациях Республики Беларусь: постановление Правления Национального
банка Республики Беларусь от 29.08.2013 №\,506 [Электронный ресурс]. --- Режим доступа:
\url{https://www.nbrb.by/legislation/accountingreporting}. --- Дата доступа: 01.10.2026.
\item Вклады для физических лиц~/ ЗАО <<Банк <<Решение>> [Электронный ресурс]. --- Режим
доступа: \url{https://rbank.by/life/deposits/}. --- Дата доступа: 01.10.2026.
\item Ньюмен, С. Создание микросервисов~/ С.~Ньюмен. --- 2-е изд. --- СПб.: Питер, 2023. ---
624~с.
\item Ричардсон, К. Микросервисы. Паттерны разработки и рефакторинга~/ К.~Ричардсон. ---
СПб.: Питер, 2019. --- 544~с.
\item Spring Framework Reference: REST Clients, HTTP Interface [Электронный ресурс]. ---
Режим доступа: \url{https://docs.spring.io/spring-framework/reference/integration/rest-clients.html}.
--- Дата доступа: 01.10.2026.
\end{enumerate}

\appendixtitle{А}{Исходный код сервисов счетов и депозитов}

% Файл сгенерирован scripts/gen_listings.py — вручную не править

\subsection*{Общая библиотека сервисов (bank-common)}

\lstinputlisting[style=xml, style=appendix, caption={bank-common/pom.xml}]{../../bank-common/pom.xml}
\lstinputlisting[style=java, style=appendix, caption={bank-common/src/main/java/by/bsuir/bank/common/HttpApis.java}]{../../bank-common/src/main/java/by/bsuir/bank/common/HttpApis.java}
\lstinputlisting[style=java, style=appendix, caption={bank-common/src/main/java/by/bsuir/bank/common/IntegrationConfig.java}]{../../bank-common/src/main/java/by/bsuir/bank/common/IntegrationConfig.java}
\lstinputlisting[style=java, style=appendix, caption={bank-common/src/main/java/by/bsuir/bank/common/ServiceUrls.java}]{../../bank-common/src/main/java/by/bsuir/bank/common/ServiceUrls.java}
\lstinputlisting[style=java, style=appendix, caption={bank-common/src/main/java/by/bsuir/bank/common/api/ApiError.java}]{../../bank-common/src/main/java/by/bsuir/bank/common/api/ApiError.java}
\lstinputlisting[style=java, style=appendix, caption={bank-common/src/main/java/by/bsuir/bank/common/api/BankException.java}]{../../bank-common/src/main/java/by/bsuir/bank/common/api/BankException.java}
\lstinputlisting[style=java, style=appendix, caption={bank-common/src/main/java/by/bsuir/bank/common/api/CommonExceptionHandler.java}]{../../bank-common/src/main/java/by/bsuir/bank/common/api/CommonExceptionHandler.java}
\lstinputlisting[style=java, style=appendix, caption={bank-common/src/main/java/by/bsuir/bank/common/client/ClientApi.java}]{../../bank-common/src/main/java/by/bsuir/bank/common/client/ClientApi.java}
\lstinputlisting[style=java, style=appendix, caption={bank-common/src/main/java/by/bsuir/bank/common/client/ClientInfo.java}]{../../bank-common/src/main/java/by/bsuir/bank/common/client/ClientInfo.java}
\lstinputlisting[style=java, style=appendix, caption={bank-common/src/main/java/by/bsuir/bank/common/ledger/AccountInfo.java}]{../../bank-common/src/main/java/by/bsuir/bank/common/ledger/AccountInfo.java}
\lstinputlisting[style=java, style=appendix, caption={bank-common/src/main/java/by/bsuir/bank/common/ledger/BankDayInfo.java}]{../../bank-common/src/main/java/by/bsuir/bank/common/ledger/BankDayInfo.java}
\lstinputlisting[style=java, style=appendix, caption={bank-common/src/main/java/by/bsuir/bank/common/ledger/Batch.java}]{../../bank-common/src/main/java/by/bsuir/bank/common/ledger/Batch.java}
\lstinputlisting[style=java, style=appendix, caption={bank-common/src/main/java/by/bsuir/bank/common/ledger/CurrencyInfo.java}]{../../bank-common/src/main/java/by/bsuir/bank/common/ledger/CurrencyInfo.java}
\lstinputlisting[style=java, style=appendix, caption={bank-common/src/main/java/by/bsuir/bank/common/ledger/DayCloseInfo.java}]{../../bank-common/src/main/java/by/bsuir/bank/common/ledger/DayCloseInfo.java}
\lstinputlisting[style=java, style=appendix, caption={bank-common/src/main/java/by/bsuir/bank/common/ledger/Entry.java}]{../../bank-common/src/main/java/by/bsuir/bank/common/ledger/Entry.java}
\lstinputlisting[style=java, style=appendix, caption={bank-common/src/main/java/by/bsuir/bank/common/ledger/LedgerApi.java}]{../../bank-common/src/main/java/by/bsuir/bank/common/ledger/LedgerApi.java}
\lstinputlisting[style=java, style=appendix, caption={bank-common/src/main/java/by/bsuir/bank/common/ledger/OpenAccount.java}]{../../bank-common/src/main/java/by/bsuir/bank/common/ledger/OpenAccount.java}
\lstinputlisting[style=java, style=appendix, caption={bank-common/src/main/java/by/bsuir/bank/common/ledger/OperationInfo.java}]{../../bank-common/src/main/java/by/bsuir/bank/common/ledger/OperationInfo.java}
\lstinputlisting[style=java, style=appendix, caption={bank-common/src/main/java/by/bsuir/bank/common/ledger/Posting.java}]{../../bank-common/src/main/java/by/bsuir/bank/common/ledger/Posting.java}
\lstinputlisting[style=java, style=appendix, caption={bank-common/src/main/java/by/bsuir/bank/common/ledger/Side.java}]{../../bank-common/src/main/java/by/bsuir/bank/common/ledger/Side.java}
\lstinputlisting[style=java, style=appendix, caption={bank-common/src/test/java/by/bsuir/bank/common/ledger/InMemoryLedger.java}]{../../bank-common/src/test/java/by/bsuir/bank/common/ledger/InMemoryLedger.java}

\subsection*{Микросервис «Счета» (account-service)}

\lstinputlisting[style=xml, style=appendix, caption={account-service/pom.xml}]{../../account-service/pom.xml}
\lstinputlisting[style=yaml, style=appendix, caption={account-service/src/main/resources/application.yml}]{../../account-service/src/main/resources/application.yml}
\lstinputlisting[style=sql, style=appendix, caption={account-service/src/main/resources/schema.sql}]{../../account-service/src/main/resources/schema.sql}
\lstinputlisting[style=sql, style=appendix, caption={account-service/src/main/resources/data.sql}]{../../account-service/src/main/resources/data.sql}
\lstinputlisting[style=java, style=appendix, caption={account-service/src/main/java/by/bsuir/bank/account/AccountServiceApplication.java}]{../../account-service/src/main/java/by/bsuir/bank/account/AccountServiceApplication.java}
\lstinputlisting[style=java, style=appendix, caption={account-service/src/main/java/by/bsuir/bank/account/BankProperties.java}]{../../account-service/src/main/java/by/bsuir/bank/account/BankProperties.java}
\lstinputlisting[style=java, style=appendix, caption={account-service/src/main/java/by/bsuir/bank/account/domain/Account.java}]{../../account-service/src/main/java/by/bsuir/bank/account/domain/Account.java}
\lstinputlisting[style=java, style=appendix, caption={account-service/src/main/java/by/bsuir/bank/account/domain/AccountNumbers.java}]{../../account-service/src/main/java/by/bsuir/bank/account/domain/AccountNumbers.java}
\lstinputlisting[style=java, style=appendix, caption={account-service/src/main/java/by/bsuir/bank/account/domain/Activity.java}]{../../account-service/src/main/java/by/bsuir/bank/account/domain/Activity.java}
\lstinputlisting[style=java, style=appendix, caption={account-service/src/main/java/by/bsuir/bank/account/domain/BankDay.java}]{../../account-service/src/main/java/by/bsuir/bank/account/domain/BankDay.java}
\lstinputlisting[style=java, style=appendix, caption={account-service/src/main/java/by/bsuir/bank/account/domain/ChartAccount.java}]{../../account-service/src/main/java/by/bsuir/bank/account/domain/ChartAccount.java}
\lstinputlisting[style=java, style=appendix, caption={account-service/src/main/java/by/bsuir/bank/account/domain/Currency.java}]{../../account-service/src/main/java/by/bsuir/bank/account/domain/Currency.java}
\lstinputlisting[style=java, style=appendix, caption={account-service/src/main/java/by/bsuir/bank/account/domain/Operation.java}]{../../account-service/src/main/java/by/bsuir/bank/account/domain/Operation.java}
\lstinputlisting[style=java, style=appendix, caption={account-service/src/main/java/by/bsuir/bank/account/domain/OperationEntry.java}]{../../account-service/src/main/java/by/bsuir/bank/account/domain/OperationEntry.java}
\lstinputlisting[style=java, style=appendix, caption={account-service/src/main/java/by/bsuir/bank/account/domain/PostingBatch.java}]{../../account-service/src/main/java/by/bsuir/bank/account/domain/PostingBatch.java}
\lstinputlisting[style=java, style=appendix, caption={account-service/src/main/java/by/bsuir/bank/account/domain/Side.java}]{../../account-service/src/main/java/by/bsuir/bank/account/domain/Side.java}
\lstinputlisting[style=java, style=appendix, caption={account-service/src/main/java/by/bsuir/bank/account/repository/AccountRepository.java}]{../../account-service/src/main/java/by/bsuir/bank/account/repository/AccountRepository.java}
\lstinputlisting[style=java, style=appendix, caption={account-service/src/main/java/by/bsuir/bank/account/repository/BankDayRepository.java}]{../../account-service/src/main/java/by/bsuir/bank/account/repository/BankDayRepository.java}
\lstinputlisting[style=java, style=appendix, caption={account-service/src/main/java/by/bsuir/bank/account/repository/ChartAccountRepository.java}]{../../account-service/src/main/java/by/bsuir/bank/account/repository/ChartAccountRepository.java}
\lstinputlisting[style=java, style=appendix, caption={account-service/src/main/java/by/bsuir/bank/account/repository/CurrencyRepository.java}]{../../account-service/src/main/java/by/bsuir/bank/account/repository/CurrencyRepository.java}
\lstinputlisting[style=java, style=appendix, caption={account-service/src/main/java/by/bsuir/bank/account/repository/OperationRepository.java}]{../../account-service/src/main/java/by/bsuir/bank/account/repository/OperationRepository.java}
\lstinputlisting[style=java, style=appendix, caption={account-service/src/main/java/by/bsuir/bank/account/repository/PostingBatchRepository.java}]{../../account-service/src/main/java/by/bsuir/bank/account/repository/PostingBatchRepository.java}
\lstinputlisting[style=java, style=appendix, caption={account-service/src/main/java/by/bsuir/bank/account/dto/AccountReport.java}]{../../account-service/src/main/java/by/bsuir/bank/account/dto/AccountReport.java}
\lstinputlisting[style=java, style=appendix, caption={account-service/src/main/java/by/bsuir/bank/account/dto/AccountView.java}]{../../account-service/src/main/java/by/bsuir/bank/account/dto/AccountView.java}
\lstinputlisting[style=java, style=appendix, caption={account-service/src/main/java/by/bsuir/bank/account/dto/BatchRequest.java}]{../../account-service/src/main/java/by/bsuir/bank/account/dto/BatchRequest.java}
\lstinputlisting[style=java, style=appendix, caption={account-service/src/main/java/by/bsuir/bank/account/dto/DayCloseResult.java}]{../../account-service/src/main/java/by/bsuir/bank/account/dto/DayCloseResult.java}
\lstinputlisting[style=java, style=appendix, caption={account-service/src/main/java/by/bsuir/bank/account/dto/EntryRequest.java}]{../../account-service/src/main/java/by/bsuir/bank/account/dto/EntryRequest.java}
\lstinputlisting[style=java, style=appendix, caption={account-service/src/main/java/by/bsuir/bank/account/dto/OpenAccountRequest.java}]{../../account-service/src/main/java/by/bsuir/bank/account/dto/OpenAccountRequest.java}
\lstinputlisting[style=java, style=appendix, caption={account-service/src/main/java/by/bsuir/bank/account/dto/OperationRequest.java}]{../../account-service/src/main/java/by/bsuir/bank/account/dto/OperationRequest.java}
\lstinputlisting[style=java, style=appendix, caption={account-service/src/main/java/by/bsuir/bank/account/dto/OperationView.java}]{../../account-service/src/main/java/by/bsuir/bank/account/dto/OperationView.java}
\lstinputlisting[style=java, style=appendix, caption={account-service/src/main/java/by/bsuir/bank/account/service/BankDayService.java}]{../../account-service/src/main/java/by/bsuir/bank/account/service/BankDayService.java}
\lstinputlisting[style=java, style=appendix, caption={account-service/src/main/java/by/bsuir/bank/account/service/DayCloseNotifier.java}]{../../account-service/src/main/java/by/bsuir/bank/account/service/DayCloseNotifier.java}
\lstinputlisting[style=java, style=appendix, caption={account-service/src/main/java/by/bsuir/bank/account/service/LedgerInitializer.java}]{../../account-service/src/main/java/by/bsuir/bank/account/service/LedgerInitializer.java}
\lstinputlisting[style=java, style=appendix, caption={account-service/src/main/java/by/bsuir/bank/account/service/LedgerService.java}]{../../account-service/src/main/java/by/bsuir/bank/account/service/LedgerService.java}
\lstinputlisting[style=java, style=appendix, caption={account-service/src/main/java/by/bsuir/bank/account/web/BankDayController.java}]{../../account-service/src/main/java/by/bsuir/bank/account/web/BankDayController.java}
\lstinputlisting[style=java, style=appendix, caption={account-service/src/main/java/by/bsuir/bank/account/web/LedgerController.java}]{../../account-service/src/main/java/by/bsuir/bank/account/web/LedgerController.java}
\lstinputlisting[style=xml, style=appendix, caption={account-service/src/main/resources/static/index.html}]{../../account-service/src/main/resources/static/index.html}
\lstinputlisting[style=js, style=appendix, caption={account-service/src/main/resources/static/app.js}]{../../account-service/src/main/resources/static/app.js}
\lstinputlisting[style=java, style=appendix, caption={account-service/src/test/java/by/bsuir/bank/account/AccountNumbersTest.java}]{../../account-service/src/test/java/by/bsuir/bank/account/AccountNumbersTest.java}
\lstinputlisting[style=java, style=appendix, caption={account-service/src/test/java/by/bsuir/bank/account/BankDayIT.java}]{../../account-service/src/test/java/by/bsuir/bank/account/BankDayIT.java}
\lstinputlisting[style=java, style=appendix, caption={account-service/src/test/java/by/bsuir/bank/account/LedgerServiceIT.java}]{../../account-service/src/test/java/by/bsuir/bank/account/LedgerServiceIT.java}

\subsection*{Микросервис «Депозиты» (deposit-service)}

\lstinputlisting[style=xml, style=appendix, caption={deposit-service/pom.xml}]{../../deposit-service/pom.xml}
\lstinputlisting[style=yaml, style=appendix, caption={deposit-service/src/main/resources/application.yml}]{../../deposit-service/src/main/resources/application.yml}
\lstinputlisting[style=sql, style=appendix, caption={deposit-service/src/main/resources/schema.sql}]{../../deposit-service/src/main/resources/schema.sql}
\lstinputlisting[style=sql, style=appendix, caption={deposit-service/src/main/resources/data.sql}]{../../deposit-service/src/main/resources/data.sql}
\lstinputlisting[style=java, style=appendix, caption={deposit-service/src/main/java/by/bsuir/bank/deposit/DepositServiceApplication.java}]{../../deposit-service/src/main/java/by/bsuir/bank/deposit/DepositServiceApplication.java}
\lstinputlisting[style=java, style=appendix, caption={deposit-service/src/main/java/by/bsuir/bank/deposit/domain/DepositContract.java}]{../../deposit-service/src/main/java/by/bsuir/bank/deposit/domain/DepositContract.java}
\lstinputlisting[style=java, style=appendix, caption={deposit-service/src/main/java/by/bsuir/bank/deposit/domain/DepositKind.java}]{../../deposit-service/src/main/java/by/bsuir/bank/deposit/domain/DepositKind.java}
\lstinputlisting[style=java, style=appendix, caption={deposit-service/src/main/java/by/bsuir/bank/deposit/domain/DepositProduct.java}]{../../deposit-service/src/main/java/by/bsuir/bank/deposit/domain/DepositProduct.java}
\lstinputlisting[style=java, style=appendix, caption={deposit-service/src/main/java/by/bsuir/bank/deposit/domain/DepositStatus.java}]{../../deposit-service/src/main/java/by/bsuir/bank/deposit/domain/DepositStatus.java}
\lstinputlisting[style=java, style=appendix, caption={deposit-service/src/main/java/by/bsuir/bank/deposit/repository/DepositContractRepository.java}]{../../deposit-service/src/main/java/by/bsuir/bank/deposit/repository/DepositContractRepository.java}
\lstinputlisting[style=java, style=appendix, caption={deposit-service/src/main/java/by/bsuir/bank/deposit/repository/DepositProductRepository.java}]{../../deposit-service/src/main/java/by/bsuir/bank/deposit/repository/DepositProductRepository.java}
\lstinputlisting[style=java, style=appendix, caption={deposit-service/src/main/java/by/bsuir/bank/deposit/dto/DepositDetails.java}]{../../deposit-service/src/main/java/by/bsuir/bank/deposit/dto/DepositDetails.java}
\lstinputlisting[style=java, style=appendix, caption={deposit-service/src/main/java/by/bsuir/bank/deposit/dto/DepositRequest.java}]{../../deposit-service/src/main/java/by/bsuir/bank/deposit/dto/DepositRequest.java}
\lstinputlisting[style=java, style=appendix, caption={deposit-service/src/main/java/by/bsuir/bank/deposit/dto/DepositView.java}]{../../deposit-service/src/main/java/by/bsuir/bank/deposit/dto/DepositView.java}
\lstinputlisting[style=java, style=appendix, caption={deposit-service/src/main/java/by/bsuir/bank/deposit/dto/Meta.java}]{../../deposit-service/src/main/java/by/bsuir/bank/deposit/dto/Meta.java}
\lstinputlisting[style=java, style=appendix, caption={deposit-service/src/main/java/by/bsuir/bank/deposit/service/DayCloseProcessor.java}]{../../deposit-service/src/main/java/by/bsuir/bank/deposit/service/DayCloseProcessor.java}
\lstinputlisting[style=java, style=appendix, caption={deposit-service/src/main/java/by/bsuir/bank/deposit/service/DepositPostings.java}]{../../deposit-service/src/main/java/by/bsuir/bank/deposit/service/DepositPostings.java}
\lstinputlisting[style=java, style=appendix, caption={deposit-service/src/main/java/by/bsuir/bank/deposit/service/DepositService.java}]{../../deposit-service/src/main/java/by/bsuir/bank/deposit/service/DepositService.java}
\lstinputlisting[style=java, style=appendix, caption={deposit-service/src/main/java/by/bsuir/bank/deposit/service/InterestCalculator.java}]{../../deposit-service/src/main/java/by/bsuir/bank/deposit/service/InterestCalculator.java}
\lstinputlisting[style=java, style=appendix, caption={deposit-service/src/main/java/by/bsuir/bank/deposit/web/DayCloseController.java}]{../../deposit-service/src/main/java/by/bsuir/bank/deposit/web/DayCloseController.java}
\lstinputlisting[style=java, style=appendix, caption={deposit-service/src/main/java/by/bsuir/bank/deposit/web/DepositController.java}]{../../deposit-service/src/main/java/by/bsuir/bank/deposit/web/DepositController.java}
\lstinputlisting[style=java, style=appendix, caption={deposit-service/src/main/java/by/bsuir/bank/deposit/web/ReferenceController.java}]{../../deposit-service/src/main/java/by/bsuir/bank/deposit/web/ReferenceController.java}
\lstinputlisting[style=xml, style=appendix, caption={deposit-service/src/main/resources/static/index.html}]{../../deposit-service/src/main/resources/static/index.html}
\lstinputlisting[style=js, style=appendix, caption={deposit-service/src/main/resources/static/app.js}]{../../deposit-service/src/main/resources/static/app.js}
\lstinputlisting[style=java, style=appendix, caption={deposit-service/src/test/java/by/bsuir/bank/deposit/DepositApiIT.java}]{../../deposit-service/src/test/java/by/bsuir/bank/deposit/DepositApiIT.java}
\lstinputlisting[style=java, style=appendix, caption={deposit-service/src/test/java/by/bsuir/bank/deposit/DepositFlowIT.java}]{../../deposit-service/src/test/java/by/bsuir/bank/deposit/DepositFlowIT.java}
\lstinputlisting[style=java, style=appendix, caption={deposit-service/src/test/java/by/bsuir/bank/deposit/InterestCalculatorTest.java}]{../../deposit-service/src/test/java/by/bsuir/bank/deposit/InterestCalculatorTest.java}
\lstinputlisting[style=java, style=appendix, caption={deposit-service/src/test/java/by/bsuir/bank/deposit/TestLedgerConfig.java}]{../../deposit-service/src/test/java/by/bsuir/bank/deposit/TestLedgerConfig.java}


\end{document}
